\chapter{Implementação da Solução}
\label{cap:implementação}

% TODO: Breve introdução ligando ao capítulo anterior — as decisões de arquitetura
% são aqui concretizadas. O capítulo organiza-se em três eixos: o desenvolvimento
% do backend (módulos, autenticação, API), a estratégia de testes, e a infraestrutura
% de produção. O microserviço de IA é mencionado na medida em que o backend o integra.

%----------------------------------------------------------------------------------------
\section{Desenvolvimento do Backend}
\label{sec:implementacao_backend}

% 1. ORGANIZAÇÃO EM MÓDULOS
%    Descrever como o NestJS estrutura a aplicação em módulos independentes:
%    AuthModule, UsersModule, InstitutionsModule, SeniorsModule, CaregiversModule,
%    ConversationsModule, MedicationsModule.
%    Cada módulo encapsula o seu controlador, serviço e repositório — nenhum módulo
%    acede diretamente ao repositório de outro; a comunicação é feita via injeção
%    de serviços.
%    Explicar como esta organização facilitou o desenvolvimento paralelo de funcionalidades
%    e reduziu o risco de regressões entre módulos.
%
% 2. AUTENTICAÇÃO E AUTORIZAÇÃO
%    Descrever o fluxo de autenticação: o utilizador envia credenciais, o AuthService
%    valida contra a hash Argon2 armazenada, e emite um JWT com o perfil (role) e o
%    identificador da instituição no payload.
%    Explicar como os guards do NestJS consomem o JWT em cada pedido: o JwtAuthGuard
%    valida a assinatura e extrai o payload; o RolesGuard verifica se o perfil do
%    utilizador tem permissão para o endpoint em questão.
%    Dar um exemplo concreto: o endpoint de criação de senior só é acessível a
%    utilizadores com o papel ADMIN da mesma instituição.
%    Mencionar a proteção contra acesso cruzado entre instituições: cada query ao
%    repositório filtra sempre pelo institutionId extraído do JWT, garantindo que
%    um administrador de uma instituição nunca acede a dados de outra.
%
% 3. MÓDULO DE ADMINISTRAÇÃO
%    Descrever as funcionalidades implementadas: criação de caregivers e seniores,
%    associação de seniores a caregivers, gestão do perfil clínico do senior (notas,
%    medicação, preferências de interação com o assistente de IA).
%    Destacar que o perfil do senior alimenta diretamente o contexto enviado ao
%    microserviço de IA — o backend é responsável por construir e enviar esse contexto
%    em cada pedido, sem que o microserviço de IA precise de consultar a base de dados.
%
% 4. VALIDAÇÃO DE DADOS E CONTRATOS DE API
%    Descrever a estratégia de validação de input com schemas partilhados entre
%    frontend e backend (definidos uma única vez, usados nos dois lados).
%    Explicar como isto garante que o contrato de API é verificado em tempo de
%    compilação TypeScript e em tempo de execução, eliminando uma categoria inteira
%    de bugs de integração.

%----------------------------------------------------------------------------------------
\section{Integração com o Serviço de Inteligência Artificial}
\label{sec:implementacao_ai}

% Descrever a integração do backend com o microserviço de IA de forma sucinta,
% focando na responsabilidade do backend neste fluxo.
%
% O QUE O BACKEND FAZ NESTE FLUXO:
%   - Autentica o pedido de voz do senior antes de o encaminhar para o microserviço.
%   - Constrói o contexto da conversa (perfil do senior, historial recente, medicação)
%     e envia-o junto com o áudio para o microserviço de IA.
%   - Persiste a conversa e o sumário estruturado devolvido pelo microserviço.
%   - Garante que o microserviço de IA nunca está exposto diretamente ao exterior —
%     todo o tráfego passa pelo backend autenticado.
%
% Não é necessário detalhar o pipeline interno do microserviço (STT → LLM → TTS)
% neste capítulo — uma descrição de alto nível é suficiente, dado que o foco do
% estágio foi a camada de backend e infraestrutura.

%----------------------------------------------------------------------------------------
\section{Estratégia de Testes}
\label{sec:testes}

% Descrever a pirâmide de testes implementada com exemplos concretos de cada nível.
% O framework adotado é o Vitest pela compatibilidade nativa com TypeScript/ESM.
%
% NÍVEL 1 — TESTES UNITÁRIOS (base da pirâmide, maior volume):
%   Cobertura da lógica de negócio dos serviços NestJS com mocks das dependências.
%   As dependências (repositórios, clientes externos) são substituídas por mocks via
%   vi.fn(), permitindo testar a lógica de negócio de forma isolada e determinística.
%   Exemplos concretos a incluir:
%     - Teste ao AuthService: verificar que a password é encriptada com Argon2 antes
%       de persistir e que um JWT é emitido com o payload correto.
%     - Teste ao SeniorService: verificar que a criação de um senior falha se já existir
%       um utilizador com o mesmo email na mesma instituição.
%     - Teste ao ConversationService: verificar que o sumário é persistido com os campos
%       corretos após o microserviço de IA devolver a resposta.
%
% NÍVEL 2 — TESTES DE CONTRATO:
%   Validação dos schemas de dados partilhados entre frontend e backend.
%   Garantem que alterações ao modelo de dados (e.g. renomear um campo) são detetadas
%   antes de chegarem a produção — o build falha se o contrato for quebrado.
%
% NÍVEL 3 — TESTES DE INTEGRAÇÃO:
%   Testes sobre os clientes dos fornecedores externos de IA, executados em ambiente
%   controlado com respostas pré-gravadas (fixtures).
%   Verificam que o cliente lida corretamente com respostas de erro, timeouts e
%   respostas parciais (streaming), sem depender da disponibilidade real dos serviços.
%
% NÍVEL 4 — TESTES DE CONFIGURAÇÃO DE DEPLOY:
%   Validação das variáveis de ambiente críticas antes de qualquer publicação.
%   O processo de deploy falha explicitamente se uma variável obrigatória estiver
%   ausente ou com formato inválido, evitando falhas silenciosas em produção.
%   Exemplo: verificar que DATABASE_URL tem o formato de connection string PostgreSQL,
%   que as chaves de API dos fornecedores externos estão presentes e não são strings
%   vazias, e que o segredo JWT tem comprimento mínimo de segurança.
%
% Incluir métricas de cobertura obtidas e análise do que ficou fora do âmbito dos testes.

%----------------------------------------------------------------------------------------
\section{Infraestrutura e Processo de Entrega}
\label{sec:deployment}

% Esta secção é um dos focos centrais do estágio. Descrever com detalhe suficiente
% para que o processo possa ser reproduzido por outro elemento da equipa.
%
% 1. CONTENTORIZAÇÃO E REDE
%    Descrever a organização dos quatro contentores Docker (frontend/Nginx, backend,
%    serviço de IA, PostgreSQL) e a rede interna que os interliga.
%    Enfatizar que apenas o contentor de proxy está ligado à rede pública — os
%    restantes comunicam exclusivamente na rede interna Docker, o que cumpre o
%    requisito RNF01 de isolamento do serviço de IA e da base de dados.
%
% 2. PROXY E TLS
%    Descrever a configuração do Nginx como reverse proxy: routing de pedidos para
%    os serviços internos com base no path, terminação TLS com certificados Let's
%    Encrypt via certbot, e headers de segurança (HSTS, X-Frame-Options, etc.).
%    Mencionar a automatização da renovação de certificados.
%
% 3. ESTRATÉGIA BLUE-GREEN
%    Descrever o modelo de dois ambientes paralelos (blue e green) e o processo de
%    atualização sem downtime:
%      a) o novo build é colocado no ambiente inativo;
%      b) os testes de configuração de deploy são executados antes da promoção;
%      c) o proxy é reconfigurado para apontar para o novo ambiente;
%      d) o ambiente anterior fica em standby para rollback imediato se necessário.
%    Explicar porque esta estratégia é especialmente relevante num serviço de cuidados
%    em que uma indisponibilidade tem impacto direto em utilizadores vulneráveis.
%
% 4. MIGRAÇÕES AUTOMÁTICAS DE BASE DE DADOS
%    Descrever como as migrações Drizzle são executadas automaticamente no arranque
%    do contentor de backend, antes de o serviço começar a aceitar pedidos.
%    Explicar as garantias desta abordagem: o código em execução é sempre compatível
%    com o esquema de dados em vigor; não é possível colocar uma versão em produção
%    com um esquema desatualizado.
%    Mencionar os cuidados necessários para migrações destrutivas (e.g. remoção de
%    colunas) em ambiente de produção com dados reais.
%
% 5. DESAFIOS ENCONTRADOS
%    - Resolução de problemas de ordem de arranque dos contentores (o backend não pode
%      iniciar antes da base de dados estar pronta) e a solução adotada.
%    - Gestão de segredos em produção: variáveis de ambiente vs. ficheiros de secrets.
%    - Configuração inicial da instância EC2: regras de firewall, user data scripts,
%      e automatização do arranque dos serviços após reinício da instância.
